\section{Enformer：从 DNA 序列预测长程调控与基因表达}

DeepConsensus 先改善“看见的序列是否正确”，Enformer 再问“这段序列会产生什么调控效果”。许多疾病相关变异位于 non-coding region，不直接改变蛋白质编码，却可能影响 transcription factor binding、enhancer--promoter interaction 与 gene expression。要预测这种效应，模型需要覆盖远超局部 motif 的上下文。

\subsection{为什么 GWAS 把问题推向非编码区}

论文页引出 Enformer，GWAS 页显示随着 cohort 增大，发现的 disease-associated loci 快速增加。Coding versus non-coding 页给出关键事实：大量关联信号位于不编码蛋白质的区域。它们可能通过调控影响基因，而不是直接改变 amino acid；因此，传统只看 coding mutation 的流程无法解释多数关联。

\lecturefigure{slide-69-gene-expression-paper.jpg}{指定工作：Effective Gene Expression Prediction from Sequence}{官方录像恢复幻灯片 V069；Stanford CS25 Lecture 17}

\lecturefigure{slide-70-gwas-growth.jpg}{GWAS 规模增长与关联位点发现}{官方录像恢复幻灯片 V070；Stanford CS25 Lecture 17}

\lecturefigure{slide-71-coding-vs-noncoding.jpg}{Coding 与 non-coding 变异的功能差异}{官方录像恢复幻灯片 V071；Stanford CS25 Lecture 17}

\begin{knowledgebox}{背景概念：GWAS、variant 与因果}
\term{GWAS} 比较群体中的 genetic variant 与表型关联；\term{variant} 是个体相对参考基因组的碱基差异；\term{association} 表示统计相关，不自动等于 causal variant。连锁不平衡、群体结构和环境因素都可能让附近多个位点共同显著，因此模型预测主要用于优先级排序与机制假设。
\end{knowledgebox}

\subsection{Enhancer、promoter 与远程作用}

Transcription regulation 页展示 promoter、enhancer、transcription factor 与 gene expression 的关系。Variant disrupts binding 页进一步说明，一个非编码碱基变化可能改变因子结合，继而改变远处基因的表达。局部卷积能识别 motif，却难以直接连接相隔数万碱基的 enhancer 与 promoter；这正是 attention 在基因组中的动机。

\lecturefigure{slide-72-transcription-regulation.jpg}{转录调控：promoter、enhancer 与 transcription factor}{官方录像恢复幻灯片 V072；Stanford CS25 Lecture 17}

\lecturefigure{slide-73-variant-disrupts-binding.jpg}{非编码变异如何破坏结合并改变表达}{官方录像恢复幻灯片 V073；Stanford CS25 Lecture 17}

\begin{warningbox}{调控预测不是疾病诊断}
模型可以预测某个 allele 可能改变 chromatin 或 expression track，但疾病表型还涉及组织、发育阶段、环境、多个基因和临床背景。变异效应分数适合筛选实验候选，不应直接翻译成个人患病概率。
\end{warningbox}

\subsection{Enformer 的输入输出与混合架构}

任务页把长 DNA sequence 作为输入，预测许多 genomic tracks；Enformer 页展示 convolution 与 Transformer 的组合。卷积先压缩局部 motif，attention 再在更长范围建立位置交互，输出头为不同细胞类型和实验测量预测信号。模型大约读取二十万碱基上下文，因此比短窗口模型更有机会连接远程调控元件。

\lecturefigure{slide-74-dna-sequence-prediction-task.jpg}{从 DNA 序列预测多种功能基因组信号}{官方录像恢复幻灯片 V074；Stanford CS25 Lecture 17}

\lecturefigure{slide-75-enformer-model.jpg}{Enformer：卷积与 Transformer 结合的长程模型}{官方录像恢复幻灯片 V075；Stanford CS25 Lecture 17}

可把模型抽象为
$$
\hat Z=f_{\theta}(X),\qquad
X\in\{A,C,G,T\}^{L},\quad
\hat Z\in\mathbb{R}^{B\times K}.
$$
其中 $X$ 是长度为 $L$ 的 DNA 序列，$B$ 是输出位置 bin 数，$K$ 是不同 assay、cell type 或 genomic track 的数量，$\hat Z$ 是预测信号。对某个变异的 effect 可用两条 allele 序列之差表示：
$$
\Delta_k=f_{\theta}(X_{\text{alt}})_k-f_{\theta}(X_{\text{ref}})_k.
$$
其中 $X_{\text{ref}}$ 与 $X_{\text{alt}}$ 只在候选 variant 处不同，$\Delta_k$ 是第 $k$ 个 track 的预测变化。这个差值是模型内的 counterfactual estimate，不是实验因果证明。

混合架构也体现了尺度分工。卷积对局部 motif 具有强归纳偏置，能够高效识别短的转录因子结合模式；池化逐步扩大感受野并降低序列长度；Transformer 再让压缩后的远距离位置交互。若一开始就对二十万碱基做全分辨率 attention，成本很高；若只堆卷积，远程 enhancer 的信号可能经过许多层才到达 promoter。Enformer 不是“attention 取代卷积”，而是把局部模式与全局调控分别交给合适组件。

\subsection{Thousands of tracks：输出不是一个标签}

Tracks 页展示模型要同时预测大量实验信号，例如不同细胞类型中的 chromatin accessibility、transcription factor binding 与 gene expression。多任务输出让共享表示吸收相关生物过程，但也增加了 label noise 与数据不均衡；某些 track 数据丰富，某些组织和人群则明显不足。

\lecturefigure{slide-76-enformer-tracks.jpg}{Enformer 预测的多细胞类型 genomic tracks}{官方录像恢复幻灯片 V076；Stanford CS25 Lecture 17}

\begin{knowledgebox}{读图：把 genomic track 当作“位置函数”}
横轴是基因组位置，纵轴是某种实验信号强度，不同颜色或行代表 assay、cell type 或模型/真实值。先看峰的位置是否对齐，再看强度是否一致，最后检查远处 enhancer 变化能否对应 promoter 附近的 expression 改变。只看总体相关系数可能掩盖关键稀有峰。
\end{knowledgebox}

\subsection{结果与可用边界}

结果页比较 promoter 与 enhancer 相关预测，显示更长上下文提升 regulatory activity prediction。Application takeaways 页把用途归纳为解释 non-coding variation、预测 gene expression 和帮助 variant prioritization。模型的科学价值在于缩小实验搜索空间：从数以万计关联位点中找出更值得做 assay 的候选。

\lecturefigure{slide-77-promoter-enhancer-results.jpg}{Enformer 对 promoter 与 enhancer activity 的预测改进}{官方录像恢复幻灯片 V077；Stanford CS25 Lecture 17}

\lecturefigure{slide-78-application-takeaways.jpg}{Enformer 的应用结论与限制}{官方录像恢复幻灯片 V078；Stanford CS25 Lecture 17}

\teachervoice{讲者先用“绝大多数 GWAS hits 在 non-coding region”建立需求，再解释远距离 enhancer--promoter interaction。这个顺序很重要：attention 不是因为它时髦，而是因为生物机制提出了超出局部 receptive field 的问题。}

\begin{importantbox}{Enformer 的正确证据链}
长 DNA 上下文提高多种 genomic track 与 expression 预测；reference/alternate allele 的差分可以为 variant 排序；实验再检验结合、表达与表型。模型把“在哪里做实验”变得更聪明，但没有取消实验，也没有把 statistical association 自动升级为 disease mechanism。
\end{importantbox}

如果要把 Enformer 用于 variant prioritization，还需要选择与疾病组织相符的 track。对心脏疾病只看任意细胞系的总体分数，可能得到生物学上不相关的候选；对群体变异还要检查训练数据是否代表目标 ancestry。更稳妥的报告会列出模型分数、相关 cell type、附近 gene、已有 eQTL/epigenomic 证据和实验可行性，让研究者知道排序依据。模型在这里扮演的是 evidence aggregator，而不是最终裁判。

\subsection{本章小结}

Enformer 将长 DNA 序列映射为多细胞类型、多 assay 的 genomic tracks，利用卷积提取局部 motif、用 Transformer 建立长程调控关系。它能改善 gene-expression 与 variant-effect 预测，为非编码变异实验排序提供依据，但关联、预测与因果仍是三层不同证据。

